\documentclass[12pt]{graduate-design-manual}
\usepackage{graphicx}
\usepackage{listings}
\usepackage{algpseudocode}
\usepackage{xcolor}
\usepackage{tocloft}
\usepackage{natbib}
\usepackage{fancyhdr}
\usepackage{hyperref}
\captionsetup{font={small}}
\tikzset{every picture/.style={font=\fontsize{10.5pt}{12pt}\selectfont}}
\renewcommand{\normalsize}{\songti\Times\fontsize{14pt}{20pt}\selectfont}

\classid{MANUAL-01}
\title{常州工学院毕业设计LaTeX模板使用手册}
\academy{计算机信息工程学院}
\field{软件工程}
\class{通用}
\id{00000000}
\author{LaTeX工作流支持组}
\Instructor{指导教师}
\InstructorTitle{教授}
\ReviewingTeacher{评阅教师}
\ReviewingTeacherTitle{教授}
\thedate{2026}{04}

\begin{document}
\makecover

\begin{originalitystatement}
本手册旨在为常州工学院毕业设计LaTeX模板提供标准、高效的使用指导。本手册的内容经过精简与重构，以原子化、金字塔结构进行组织。
\end{originalitystatement}

\begin{cnabstract}{常州工学院毕业设计LaTeX模板使用手册}{LaTeX；工作流；排版指南}
    本文档为《常州工学院毕业设计LaTeX模板使用手册》，主要面向需要使用LaTeX进行毕业设计论文排版的学生和指导教师。为了解决过去模板使用中遇到的环境配置复杂、编译报错难排查等问题，本手册采用了金字塔原理进行内容结构化重构，剔除冗余解释，直接提供高频决策路径与排错清单。主要内容涵盖了从零开始的环境配置（包括本地编译与Overleaf线上编译两种流派）、模板定制修改的核心指令，以及常见报错的快速定位与处理方法。
    
    本手册的最终目的是极大地降低用户的学习和试错成本，实现“开箱即用、一键编译”的高效文档写作体验。
\end{cnabstract}

\contents
\pagenumbering{arabic}
\mainbody

\section{引言：工作流决策树}
\subsection{选路判定}
\begin{enumerate}
    \item 第一次接触 LaTeX 且需要免费本地编译环境：\textbf{路径 A}（新手本地化）
    \item 第一次接触 LaTeX 且需要多人协作：\textbf{路径 B}（Overleaf线上协作）
    \item 已有完整的本地编译环境，仅需使用特定模板：\textbf{路径 C}（本地模板实践）
    \item 进行常州工学院（CZU）毕业设计写作：\textbf{路径 D}（CZU 毕设专线）
\end{enumerate}

\subsection{路径执行流}
\begin{description}
    \item[A. 新手本地化] 安装 TeX Live $\rightarrow$ 安装 VS Code 及其 LaTeX Workshop 插件 $\rightarrow$ 跑通最小示例。
    \item[B. 新手协作版] 注册 Overleaf $\rightarrow$ 创建 Blank 项目或上传 ZIP 源码包 $\rightarrow$ 快捷键 \texttt{Cmd+Enter} 编译 $\rightarrow$ 若需支持中文，在导言区添加 \texttt{\textbackslash usepackage\{ctex\}}。
    \item[C. 本地模板实践] 标准编译链推荐：\texttt{xelatex $\rightarrow$ bibtex $\rightarrow$ xelatex $\rightarrow$ xelatex}。
    \item[D. CZU 毕设] 获取模板（\texttt{git clone}） $\rightarrow$ 修改 \texttt{main.tex} 基本信息 $\rightarrow$ 在 \texttt{chapters/} 目录下编写具体内容并通过 \texttt{\textbackslash input} 引入 $\rightarrow$ 执行编译（使用本地 xelatex 或上传 Overleaf）。
\end{description}


\section{环境配置与安装}
\subsection{方案 1：本地工作流 (TeX Live + VS Code)}

\subsubsection{编译环境安装}
\begin{itemize}
    \item \textbf{Mac OS}: 终端执行 \texttt{brew install {-}{-}cask mactex-no-gui}
    \item \textbf{Windows}: 官网下载 \texttt{install-tl-windows.exe}，执行完整安装。
    \item \textbf{Ubuntu}: 终端执行 \texttt{sudo apt install texlive-full}
\end{itemize}

\subsubsection{编辑器配置}
推荐使用 VS Code 配合 LaTeX Workshop 插件。
\begin{itemize}
    \item \textbf{VS Code 设置 (\texttt{.vscode/settings.json})}: 
    配置 recipes 编译链为 \texttt{xelatex $\rightarrow$ bibtex $\rightarrow$ xelatex*2}。为了节省性能，建议关闭自动编译，设置 \texttt{"latex-workshop.latex.autoBuild.run": "never"}。
    \item \textbf{快捷键配置 (\texttt{keybindings.json})}: 
    建议将 \texttt{Cmd+S} (或 Ctrl+S) 绑定为 \texttt{latex-workshop.build} 命令。
\end{itemize}

\subsubsection{核心编译指令}
标准的编译指令为 \texttt{xelatex main.tex}。复杂文档包含参考文献时需配合 bibtex。

\subsection{方案 2：云端工作流 (Overleaf)}
适合不方便配置本地环境或需要多人协作的用户。
\begin{enumerate}
    \item 注册并登录 Overleaf。
    \item 点击 New Project。
    \item 编写代码后按 \texttt{Cmd+Enter}（或 Ctrl+Enter）进行编译。
    \item \textbf{中文支持}: 需在文档导言区加入 \texttt{\textbackslash usepackage[UTF8]\{ctex\}}。
\end{enumerate}


\section{模板核心使用指令}
\subsection{核心原则}
修改前必须编译验证；每修改一处即编译一次；重要修改前先备份。

\subsection{模板目录与风险提示}
\begin{itemize}
    \item \texttt{*.cls}：全局格式定义文件，风险最高，不建议新手随意修改。
    \item \texttt{main.tex}：入口组装与基本信息配置，风险较低。
    \item \texttt{chapters/*.tex}：论文正文内容存放位置，风险最低。
\end{itemize}

\subsection{常见修改场景}
\begin{enumerate}
    \item \textbf{仅改内容}：直接在 \texttt{chapters/*.tex} 中编写。切勿修改 \texttt{\textbackslash makechaptertitle} 或 \texttt{\textbackslash ifthesis}。在 \texttt{main.tex} 中填写论文信息，如 \texttt{\textbackslash title\{\}}, \texttt{\textbackslash author\{\}} 等。
    \item \textbf{修改格式}（需修改 \texttt{*.cls} 文件）：
    \begin{itemize}
        \item 标题样式：搜索 \texttt{\textbackslash chapter}
        \item 行距：搜索 \texttt{\textbackslash linespread}
        \item 页眉页脚：搜索 \texttt{fancyhdr}
        \item 字号调整：搜索 \texttt{\textbackslash zihao}（注意 \texttt{-4} 代表小四号字）
    \end{itemize}
    \item \textbf{增删章节}：
    \begin{itemize}
        \item 增加：在 \texttt{chapters/} 新建 tex 文件，在 \texttt{main.tex} 中添加 \texttt{\textbackslash input\{chapters/xx\}}。
        \item 删除：直接注释掉对应的 \texttt{\textbackslash input\{\}} 代码行。
    \end{itemize}
    \item \textbf{参考文献样式}：
    模板通常使用国标，如果需要更换，可修改宏包参数：\\
    \texttt{\textbackslash usepackage[backend=biber,style=gb7714-2015]\{biblatex\}}
\end{enumerate}

\subsection{高频语法速查}
\begin{itemize}
    \item \textbf{文本}：粗体 \texttt{\textbackslash textbf\{\}}，斜体 \texttt{\textbackslash textit\{\}}
    \item \textbf{层级}：章 \texttt{\textbackslash chapter\{\}}，节 \texttt{\textbackslash section\{\}}，小节 \texttt{\textbackslash subsection\{\}}
    \item \textbf{公式}：行内公式 \texttt{\$ \$}，块级公式 \texttt{\textbackslash begin\{equation\} \textbackslash end\{equation\}}
    \item \textbf{图片}：\texttt{\textbackslash includegraphics[width=0.8\textbackslash textwidth]\{file.png\}}
    \item \textbf{引用}：\texttt{\textbackslash cite\{\}} 配合底部 \texttt{\textbackslash printbibliography} 或 \texttt{\textbackslash bibliography} 使用。
\end{itemize}

\subsection{常用拓展宏包}
\begin{itemize}
    \item 代码高亮：\texttt{listings}
    \item 交叉引用：\texttt{cleveref}
    \item 复杂绘图：\texttt{tikz}
\end{itemize}


\section{常见问题排错速查表}
编译 LaTeX 经常会遇到报错，以下是常见错误类型及排查手段的速查表：

\begin{table}[htbp]
\centering
\caption{LaTeX 排错速查表}
\begin{tabular}{lp{10cm}}
\hline
\textbf{错误提示} & \textbf{原因与解决办法} \\
\hline
\texttt{Undefined control sequence} & 拼写错误或未引入对应的宏包。请检查命令拼写，或在导言区添加所需的 \texttt{\textbackslash usepackage\{\}}。 \\
\texttt{File not found} & 路径或文件名错误。检查包含文件或图片的路径，确保使用正斜杠 \texttt{/} 而非反斜杠。 \\
\texttt{Missing \$ inserted} & 数学公式环境错误。检查公式中是否混入了普通文本，或者行内公式遗漏了 \texttt{\$} 符号。 \\
\texttt{Package xxx Warning} & 宏包警告，不影响最终生成 PDF，通常可以忽略。若有强迫症可根据提示调整参数。 \\
\hline
\end{tabular}
\end{table}

\subsection{通用自救指南}
\begin{enumerate}
    \item 仔细阅读终端输出的 \texttt{Error} 报错行，通常会指明报错的具体行号。
    \item 若无法定位问题，尝试将近期新增的代码注释掉，重新编译，以二分法定位错误源。
    \item 若更改了核心结构或 \texttt{.bib} 参考文献，务必清理所有的辅助文件（如 \texttt{.aux}, \texttt{.bbl}, \texttt{.log} 等）后再次完整编译。
\end{enumerate}


\end{document}
